\documentclass[../main.tex]{subfiles}
\begin{document}

\chapter{Problem Statement}
\label{ch:problem-statement}

% Owner: TBD
% Per COS700 lecture (Research Methods, Lecture 8): describe the problem
% in as much detail as possible, specify main objectives and
% sub-objectives, define scope/limitations, and state assumptions.

\section{Problem Description}
\label{sec:problem-description}
% TODO: state precisely what XAI fragility means for this project --
% explanations of cancer drug-response models that are unfaithful,
% unstable under small perturbations, or not robust across GA/NSGA-II
% runs -- and why this is a problem worth solving.

\section{Objectives}
\label{sec:problem-objectives}
% TODO: main objective + sub-objectives, phrased as "to demonstrate...",
% "to evaluate...", "to determine...". Pull directly from the proposal's
% aims/objectives, tightened to reflect what was actually implemented.

\section{Scope and Limitations}
\label{sec:problem-scope}
% TODO: explicitly state what this work does and does not deal with,
% e.g. "This work does not consider ...". Reuse proposal's scope
% section if one exists.

\section{Assumptions}
\label{sec:problem-assumptions}
% TODO: any assumptions the reader should be aware of (e.g. ANN is
% frozen and queried pedagogically, dataset coverage assumptions).

\end{document}
